\section{Enfoque metodológico}
\label{sec:enfoque}

\subsection{Paradigma y diseño de la investigación}

El proyecto se inscribe epistemológicamente en el paradigma positivista y adopta
un diseño metodológico de carácter \textbf{cuantitativo, experimental y basado
en simulación estocástica computacional}, fundamentado en la modelación dinámica
en espacio de estados, la estimación óptima recursiva y el método de Monte Carlo
\parencite{kalman1960,simon2006}.

Esta elección no es una preferencia estilística sino una consecuencia de la
pregunta de investigación. La pregunta exige cuantificar y minimizar una métrica
de error sobre variables físicas continuas, y determinar si la mejora observada
es atribuible al diseño del módulo o al azar. Un diseño cualitativo no podría
responderla, y un diseño observacional sobre datos de planta —además de
inviable por las restricciones de confidencialidad— carecería de la verdad de
referencia contra la cual medir el error, puesto que en una faena real el valor
verdadero del volumen interno es justamente lo desconocido.

La simulación resuelve ese problema epistemológico de fondo: al generar
sintéticamente la trayectoria real del sistema, se dispone de una verdad de
referencia exacta contra la cual contrastar tanto la lectura degradada como la
estimación del módulo. Esta es la condición que hace posible calcular un error
verdadero y no un residuo aparente.

\subsection{Tipo de datos y mitigación de la restricción de confidencialidad}

Debido a las estrictas políticas de confidencialidad y a los acuerdos de no
divulgación imperantes en la industria minera, el acceso a datos SCADA
históricos en crudo se encuentra restringido. Para abordar esta limitación se
utilizan \textbf{datos sintéticos estocásticos paramétricos}.

El procedimiento consiste en desarrollar un modelo de la planta física —el
\emph{truth model}— basado en las leyes de conservación de flujo, al cual se
inyecta artificialmente ruido blanco gaussiano de proceso y de medición. Las
matrices de covarianza $\mathbf{Q}$ y $\mathbf{R}$ se parametrizan dentro de
rangos operativos consistentes con la literatura técnica y con la física del
activo, según se detalla y justifica en la sección~\ref{sec:unidad}.

Conviene declarar el estatus metodológico de esta decisión con precisión. Los
datos sintéticos no sustituyen a los datos reales ni pretenden hacerlo: permiten
evaluar la \emph{capacidad del estimador} bajo condiciones controladas y
conocidas, lo que constituye una validación de laboratorio. La transferencia a
un entorno operativo requiere calibración empírica posterior, restricción que se
declara de manera explícita en las secciones~\ref{sec:etica}
y~\ref{sec:discusion}.

\subsection{Técnicas de análisis computacional}

La técnica central de análisis corresponde al \textbf{Filtro de Kalman Lineal
Discreto}. Para garantizar la reproducibilidad científica, el entorno de pruebas
utiliza Python~3 junto con las bibliotecas estandarizadas NumPy y SciPy, con
semilla aleatoria de inicialización fija y declarada.

Para validar la robustez del estimador frente a distintos niveles de degradación
instrumental se emplea el \textbf{método de simulaciones de Monte Carlo},
ejecutando $M = 50$ realizaciones independientes de un horizonte de $N = 1000$
pasos cada una. Este diseño asegura que el desempeño reportado no sea producto
de fluctuaciones estocásticas aisladas de una corrida particular, y provee la
muestra de pares sobre la cual se aplica la inferencia estadística descrita en
la sección~\ref{sec:analisis}.

\subsection{Justificación tecnológica de las herramientas}

La selección de Python con NumPy y SciPy responde a dos criterios de ingeniería.

El primero es la \textbf{eficiencia matemática nativa}. El Filtro de Kalman
requiere la resolución recurrente de álgebra lineal densa, incluyendo
multiplicación e inversión matricial en cada ciclo de tiempo discretizado. NumPy
y SciPy encapsulan subrutinas optimizadas de bajo nivel que minimizan el costo
computacional, lo que sostiene la afirmación de compatibilidad con la operación
en línea de un gemelo digital.

El segundo es el \textbf{estándar de la industria y la sostenibilidad}. Python
constituye el estándar contemporáneo para el desarrollo de \emph{pipelines} de
inteligencia artificial y ciencia de datos. El uso de bibliotecas de código
abierto elimina los costos asociados a licencias de software propietario y
garantiza la reproducibilidad científica, la portabilidad y la sostenibilidad
técnica del código a largo plazo, condiciones necesarias para que un tercero
pueda verificar los resultados reportados en este informe.
